\chapter{Discussion and future work}
\label{chap:conclusion}

Football clubs and analysts routinely ask which players resemble each other in style, yet
answer that question from memory and reputation rather than from a stated, checkable definition
of similarity. This thesis set out to replace that informal judgement with a hierarchical
Bayesian model of player-season performance capable of recovering a small number of
interpretable stylistic dimensions and of quantifying how similar any two players are, together
with the uncertainty in that judgement.

The evidence assembled across the preceding chapters supports that aim on its own terms.
Exploratory analysis established that La Liga's 48 standardised performance features are
neither redundant nor irreducibly high-dimensional: principal component analysis, factor
analysis, and multidimensional scaling all pointed to genuine, moderate-dimensional structure,
without settling on a single preferred representation. The additive hierarchical model
developed afterward resolved that ambiguity by treating player style as a shared, low-rank
component of player identity, separate from season-level drift and from residual variation
whose heavy tails and minutes-dependence needed explicit modelling rather than being averaged
away. The results chapter showed that this structure is not a modelling convenience: the fitted
covariance reproduces the observed cross-feature correlations almost exactly, and three
canonical axes, named directly from their own loadings as Defensive engagement, Technical
orientation, and Possession retention, organise the entire outfield population into recognisable
positional archetypes without positional labels ever entering the model.

Two results connect the statistical machinery back to the practical question the thesis opened
with. First, the same covariance structure that identifies the three axes also defines a single,
bounded similarity score, so that how similar two players are has one consistent answer
regardless of which pair is asked about, rather than a different ad hoc comparison each time.
Second, the four case studies show that this score behaves the way a useful measure should: it
separates genuinely rare profiles (Messi, whose closest match is still only $45\%$ similar) from
tightly clustered ones (Ramos and Piqué, at $93\%$), and, in Stuani's case, it returns neighbours
who share his style but not, by reputation, his end product.

That last result is informative in its own right: a style-only measure should not be expected to
recover output, and showing exactly where it stops is more useful than a measure that quietly
conflates the two. It is also the most direct link to the four directions taken up below, each
building on the representation established here rather than replacing it.

Beyond its use as a scouting tool, the contribution of this thesis is methodological: it
demonstrates that a fully Bayesian, identifiable factor model can be fit to a real, moderately
sized sports dataset, diagnosed rigorously, and translated into quantities, canonical axes, a
similarity score, positional archetypes, that a football analyst, not only a statistician, can
read directly. That combination of statistical accountability and domain interpretability is
what a purely descriptive or a purely predictive approach to player data cannot offer on its
own.

Four directions were considered during this work and are reported here as scoped future work rather than as
contributions of the present thesis.

\textbf{A performance score alongside similarity.} \S\ref{sec:results-similarity} deliberately measures style
only: two players can share a high similarity score while producing very differently, Stuani's case study
being the clearest example. A genuinely useful companion is a performance score, defined carefully enough to
survive scrutiny rather than a raw goals count (which conflates opportunity, minutes, and finishing quality)
or any single loose per90 rate. Paired with the similarity score, it would let a shortlist be built in two
steps instead of one: filter to players who play like the target, then rank that shortlist by this measure to
surface the best of them, rather than asking one number to do both jobs.

\textbf{Uncertainty beyond posterior means.} The model is fully Bayesian throughout, but two summaries still
reduce a full posterior to a single number: the archetype clustering of \S\ref{sec:results-stylespace} assigns
one hard $k$-means cluster per player rather than a soft, probability-weighted membership, and the similarity
scores of \S\ref{sec:results-similarity} report one distance per pair rather than a credible interval on it.
The first is a direct extension of the uncertainty already reported for canonical scores; the second is
harder, since a pairwise distance's uncertainty depends on the joint draw-by-draw behaviour of both players
together, not either one's marginal interval, and was not attempted here.

\textbf{Distance-based latent-space models.} Chapter~\ref{chap:methodology}'s factor-analytic model represents
style through a shared covariance structure; a complementary alternative represents it directly through
pairwise dissimilarities via Bayesian multidimensional scaling \cite{oh2001bayesian}, with hierarchical
extensions for repeated player--season observations \cite{yanchenko2020hierarchical}. Chapter~\ref{chap:background}'s
classical MDS is a first, non-probabilistic step in that direction; a fully Bayesian version would attach
posterior uncertainty directly to each player's position rather than obtaining it indirectly through canonical
factor scores.

\textbf{Team-season effects and tactical fit.} No team identifier exists anywhere in the current data
pipeline, so a genuine player-vs-system comparison is not buildable without new data engineering. A
team-season effect $c_{t(i,j),j}$ would separate manager and tactical-system effects from stable player
identity and would be the natural foundation for a real tactical-fit score; the closest available proxy today
is simply a player's archetype cluster (\S\ref{sec:results-stylespace}).
